\section{Harmony lemma}\label{sec:harmony}

This section proves the main result of the paper: the reductions of a
process $\conf{\rho}{P}$ (\cref{fig:sts}) and the silent transitions of the
composition $\rho \lpar P$ of its memory with its program (\cref{sec:lts})
coincide, up to structural congruence, and with the same choice.

\begin{theorem}[Harmony]\label{lem:harmony}
For all $\rho, P, r, \rho', Q$,
\[
\conf{\rho}{P} \red_{r} \conf{\rho'}{Q}
\iff
\rho \lpar P \step{\tau}{r} \rho' \lpar Q' \text{ for some } Q' \equiv Q.
\]
\end{theorem}

The rest of the section is the proof. It has two ingredients: the shapes of the transitions of a program
(\cref{lem:shapes}), and the fact that the structural congruence is a
bisimulation (\cref{lem:cgr-bisim}).

\begin{lemma}[Shapes of the transitions]\label{lem:shapes}
In the LTS of programs,
\begin{enumerate}
\item if $P \step{c!v}{r} P'$ then $r = t$ and
  $P \equiv (\tagged{t}{c!e} \parallel R) \setminus \tilde d$,
  $P' \equiv (\mathbf{0} \parallel R) \setminus \tilde d$ for some $t, e, R, \tilde d$
  with $e \Downarrow v$ and $c \notin \tilde d$;
\item if $P \step{c?v}{r} P'$ then $r = t$ and
  $P \equiv ((\tagged{t}{c?x}.P_1 + K) \parallel R) \setminus \tilde d$,
  $P' \equiv (P_1[v/x] \parallel R) \setminus \tilde d$ for some
  $t, x, P_1, K, R, \tilde d$ with $c \notin \tilde d$;
\item for the quantum actions, for some $t, P_1, K, R, \tilde d$ (and $x$
  for a measurement), with $r = t$ in the three cases:
  \begin{itemize}
  \item if $P \step{\req{U}}{r} P'$ then
    $P \equiv ((\tagged{t}{U(\tilde q)}.P_1 + K) \parallel R) \setminus \tilde d$ and
    $P' \equiv (P_1 \parallel R) \setminus \tilde d$;
  \item if $P \step{\initp}{r} P'$ then
    $P \equiv ((\tagged{t}{\initp}.P_1 + K) \parallel R) \setminus \tilde d$ and
    $P' \equiv (P_1 \parallel R) \setminus \tilde d$;
  \item if $P \step{\mrecv{m}}{r} P'$ then
    $P \equiv ((\tagged{t}{\M(\tilde q \rhd x)}.P_1 + K) \parallel R) \setminus \tilde d$ and
    $P' \equiv (P_1[m/x] \parallel R) \setminus \tilde d$;
  \end{itemize}
\end{enumerate}
\end{lemma}
\begin{proof}
  We proceed by structural induction on the rule of the LTS in the hypothesis.
\end{proof}

\begin{lemma}[Congruence is a bisimulation]\label{lem:cgr-bisim}
If $\conf{\rho}{Q} \equiv \conf{\rho}{P}$ and $\conf{\rho}{P} \step{\lambda}{r} \conf{\rho'}{P'}$ in the LTS of programs, then
$\conf{\rho}{Q} \step{\lambda}{r} \conf{\rho'}{Q'}$ for some $\conf{\rho}{Q'} \equiv \conf{\rho'}{P'}$.
\end{lemma}
\begin{proof}
  We proceed by structural induction on $\conf{\rho}{Q} \equiv \conf{\rho}{P}$.
\end{proof}

\begin{lemma}[Reductions are internal transitions]\label{lem:red-int}
If $\conf{\rho}{P} \red_{r} \conf{\rho'}{Q}$ then $P \step{\lambda}{r} Q'$ and
$\rho \eff{\lambda} \rho'$ for some internal label $\lambda$ and some
$Q' \equiv Q$.
\end{lemma}

\begin{proof}
  We proceed by structural induction on $\conf{\rho}{P} \red_{r} \conf{\rho'}{Q}$.

  We conclude directly by the rules of the LTS in all cases except for structural congruence rule.
  For structural congruence rule, we use \cref{lem:cgr-bisim} to conclude.
\end{proof}

\begin{lemma}[Internal transitions are reductions]\label{lem:int-red}
If $\conf{\rho}{P} \step{\tau}{r} \conf{\rho'}{Q}$, then $\conf{\rho}{P} \red_{r} \conf{\rho'}{Q}$.
\end{lemma}

\begin{proof}
  The proof goes by structural induction on $\conf{\rho}{P} \step{\tau}{r} \conf{\rho'}{Q}$, and we use \cref{lem:shapes} to conclude.
\end{proof}